
\documentclass{article}
\usepackage{colortbl}
\usepackage{makecell}
\usepackage{multirow}
\usepackage{supertabular}

\begin{document}

\newcounter{utterance}

\twocolumn

{ \footnotesize  \setcounter{utterance}{1}
\setlength{\tabcolsep}{0pt}
\begin{supertabular}{c@{$\;$}|p{.15\linewidth}@{}p{.15\linewidth}p{.15\linewidth}p{.15\linewidth}p{.15\linewidth}p{.15\linewidth}}

    \# & $\;$A & \multicolumn{4}{c}{Game Master} & $\;\:$B\\
    \hline 

    \theutterance \stepcounter{utterance}  

    & & \multicolumn{4}{p{0.6\linewidth}}{\cellcolor[rgb]{0.9,0.9,0.9}{%
	\makecell[{{p{\linewidth}}}]{% 
	  \tt {\tiny [A$\langle$GM]}  
	 State giocando a un gioco di negoziazione cooperativa in cui dovete accordarvi su come dividere una serie di oggetti tra voi e un altro giocatore.\\ \tt \\ \tt Regole:\\ \tt (a) A Lei e all'altro giocatore viene dato un insieme di oggetti. Ognuno di voi riceve anche una funzione di valore segreta, che rappresenta il valore di ciascun tipo di oggetto.\\ \tt (b) Si scambiano messaggi con l'altro giocatore per concordare chi si aggiudica gli oggetti. Potete inviare un massimo di 5 messaggi ciascuno, oppure terminare in anticipo facendo una proposta segreta in qualsiasi momento.\\ \tt (c) A ciascuno di voi viene chiesto di inviare una proposta segreta indicando gli oggetti che desiderate, formattata tra parentesi quadre come segue: "[Proposta: <numero> <nome oggetto>, <numero> <nome oggetto>, <...>]".\\ \tt (d) Se le vostre proposte sono complementari, cioè ci sono abbastanza oggetti per soddisfare entrambe le proposte, a ciascun giocatore viene assegnato un punteggio basato sulla somma dei valori degli oggetti ricevuti. In caso contrario, entrambi ricevono zero punti.\\ \tt (e) L'obiettivo è massimizzare la somma del proprio punteggio e di quello dell'altro giocatore.\\ \tt \\ \tt Cominciamo.\\ \tt \\ \tt L'insieme degli oggetti disponibili è:\\ \tt \\ \tt 1 telefono, 2 specchi, 2 frutti\\ \tt \\ \tt La funzione valore segreta è:\\ \tt \\ \tt telefono: 4, specchio: 2, frutta: 1\\ \tt \\ \tt IMPORTANTE: i vostri messaggi, a meno che non si tratti di una proposta segreta, vengono trasmessi direttamente all'altro giocatore, quindi non includete alcuna risposta alle regole o testo di annuncio del vostro messaggio. Per fare una proposta segreta, utilizzate il formato indicato. Non utilizzare delle parentesi quadre quando si comunica all'altro giocatore, altrimenti verrà interpretata come una proposta segreta.\\ \tt  
	  } 
	   } 
	   } 
	 & \\ 
 

    \theutterance \stepcounter{utterance}  

    & \multicolumn{4}{p{0.6\linewidth}}{\cellcolor[rgb]{0.8,1,0.9}{%
	\makecell[{{p{\linewidth}}}]{% 
	  \tt {\tiny [A$\rangle$GM]}  
	 Ciao! Ecco come valuto gli oggetti: telefono vale molto (4), specchio medio (2), frutta poco (1). Mi interessa massimizzare la somma totale.\\ \tt \\ \tt Propongo di assegnare il telefono a me (dato l’alto valore che gli do) e poi distribuire specchi e frutta in base ai tuoi valori. Puoi dirmi quanto valgono per te telefono, specchio e frutta? Se preferisci, indica quali oggetti desideri e quanto ci tieni a ciascuno, così troviamo una divisione efficiente. 
	  } 
	   } 
	   } 
	 & & \\ 
 

    \theutterance \stepcounter{utterance}  

    & & & \multicolumn{2}{p{0.3\linewidth}}{\cellcolor[rgb]{0.95,0.95,0.95}{%
	\makecell[{{p{\linewidth}}}]{% 
	  \tt {\tiny [GM$|$GM]}  
	 continue 
	  } 
	   } 
	   } 
	 & & \\ 
 

    \theutterance \stepcounter{utterance}  

    & & & \multicolumn{2}{p{0.3\linewidth}}{\cellcolor[rgb]{0.95,0.95,0.95}{%
	\makecell[{{p{\linewidth}}}]{% 
	  \tt {\tiny [GM$|$GM]}  
	 Ciao! Ecco come valuto gli oggetti: telefono vale molto (4), specchio medio (2), frutta poco (1). Mi interessa massimizzare la somma totale.\\ \tt \\ \tt Propongo di assegnare il telefono a me (dato l’alto valore che gli do) e poi distribuire specchi e frutta in base ai tuoi valori. Puoi dirmi quanto valgono per te telefono, specchio e frutta? Se preferisci, indica quali oggetti desideri e quanto ci tieni a ciascuno, così troviamo una divisione efficiente. 
	  } 
	   } 
	   } 
	 & & \\ 
 

    \theutterance \stepcounter{utterance}  

    & & \multicolumn{4}{p{0.6\linewidth}}{\cellcolor[rgb]{0.9,0.9,0.9}{%
	\makecell[{{p{\linewidth}}}]{% 
	  \tt {\tiny [GM$\rangle$B]}  
	 State giocando a un gioco di negoziazione cooperativa in cui dovete accordarvi su come dividere una serie di oggetti tra voi e un altro giocatore.\\ \tt \\ \tt Regole:\\ \tt (a) A Lei e all'altro giocatore viene dato un insieme di oggetti. Ognuno di voi riceve anche una funzione di valore segreta, che rappresenta il valore di ciascun tipo di oggetto.\\ \tt (b) Si scambiano messaggi con l'altro giocatore per concordare chi si aggiudica gli oggetti. Potete inviare un massimo di 5 messaggi ciascuno, oppure terminare in anticipo facendo una proposta segreta in qualsiasi momento.\\ \tt (c) A ciascuno di voi viene chiesto di inviare una proposta segreta indicando gli oggetti che desiderate, formattata tra parentesi quadre come segue: "[Proposta: <numero> <nome oggetto>, <numero> <nome oggetto>, <...>]".\\ \tt (d) Se le vostre proposte sono complementari, cioè ci sono abbastanza oggetti per soddisfare entrambe le proposte, a ciascun giocatore viene assegnato un punteggio basato sulla somma dei valori degli oggetti ricevuti. In caso contrario, entrambi ricevono zero punti.\\ \tt (e) L'obiettivo è massimizzare la somma del proprio punteggio e di quello dell'altro giocatore.\\ \tt \\ \tt Cominciamo.\\ \tt \\ \tt L'insieme degli oggetti disponibili è:\\ \tt \\ \tt 1 telefono, 2 specchi, 2 frutti\\ \tt \\ \tt La funzione valore segreta è:\\ \tt \\ \tt telefono: 6, specchio: 0, frutta: 2\\ \tt \\ \tt IMPORTANTE: i vostri messaggi, a meno che non si tratti di una proposta segreta, vengono trasmessi direttamente all'altro giocatore, quindi non includete alcuna risposta alle regole o testo di annuncio del vostro messaggio. Per fare una proposta segreta, utilizzate il formato indicato. Non utilizzare delle parentesi quadre quando si comunica all'altro giocatore, altrimenti verrà interpretata come una proposta segreta.\\ \tt \\ \tt \\ \tt Ciao! Ecco come valuto gli oggetti: telefono vale molto (4), specchio medio (2), frutta poco (1). Mi interessa massimizzare la somma totale.\\ \tt \\ \tt Propongo di assegnare il telefono a me (dato l’alto valore che gli do) e poi distribuire specchi e frutta in base ai tuoi valori. Puoi dirmi quanto valgono per te telefono, specchio e frutta? Se preferisci, indica quali oggetti desideri e quanto ci tieni a ciascuno, così troviamo una divisione efficiente. 
	  } 
	   } 
	   } 
	 & \\ 
 

    \theutterance \stepcounter{utterance}  

    & & & \multicolumn{4}{p{0.6\linewidth}}{\cellcolor[rgb]{1,0.85,0.72}{%
	\makecell[{{p{\linewidth}}}]{% 
	  \tt {\tiny [GM$\langle$B]}  
	 Grazie per la trasparenza. Ecco i miei valori: telefono 6, specchio 0, frutta 2 ciascuno.\\ \tt \\ \tt Dato che entrambi valutiamo molto il telefono, ma io lo valuto 6 contro il tuo 4, l’allocazione efficiente per massimizzare la somma totale è:\\ \tt - Telefono a me (valore totale 6 per me vs 4 per te)\\ \tt - Specchi a te (valgono 0 per me e 2 ciascuno per te)\\ \tt - Frutti possiamo dividerli: io ne prendo almeno 1 (vale 2 per me), tu l’altro (vale 1 per te). Somma: 6 + 4 + 2 + 1 = 13.\\ \tt \\ \tt Se sei d’accordo, posso inviare subito una proposta segreta con: 1 telefono, 1 frutto, e lascio a te 2 specchi e 1 frutto. Fammi sapere se preferisci prendere entrambi i frutti tu (per me va bene rinunciare a 2 punti se questo facilita l’accordo), oppure confermi la divisione 1-1 dei frutti. 
	  } 
	   } 
	   } 
	  \\ 
 

    \theutterance \stepcounter{utterance}  

    & & & \multicolumn{2}{p{0.3\linewidth}}{\cellcolor[rgb]{0.95,0.95,0.95}{%
	\makecell[{{p{\linewidth}}}]{% 
	  \tt {\tiny [GM$|$GM]}  
	 continue 
	  } 
	   } 
	   } 
	 & & \\ 
 

    \theutterance \stepcounter{utterance}  

    & & & \multicolumn{2}{p{0.3\linewidth}}{\cellcolor[rgb]{0.95,0.95,0.95}{%
	\makecell[{{p{\linewidth}}}]{% 
	  \tt {\tiny [GM$|$GM]}  
	 Grazie per la trasparenza. Ecco i miei valori: telefono 6, specchio 0, frutta 2 ciascuno.\\ \tt \\ \tt Dato che entrambi valutiamo molto il telefono, ma io lo valuto 6 contro il tuo 4, l’allocazione efficiente per massimizzare la somma totale è:\\ \tt - Telefono a me (valore totale 6 per me vs 4 per te)\\ \tt - Specchi a te (valgono 0 per me e 2 ciascuno per te)\\ \tt - Frutti possiamo dividerli: io ne prendo almeno 1 (vale 2 per me), tu l’altro (vale 1 per te). Somma: 6 + 4 + 2 + 1 = 13.\\ \tt \\ \tt Se sei d’accordo, posso inviare subito una proposta segreta con: 1 telefono, 1 frutto, e lascio a te 2 specchi e 1 frutto. Fammi sapere se preferisci prendere entrambi i frutti tu (per me va bene rinunciare a 2 punti se questo facilita l’accordo), oppure confermi la divisione 1-1 dei frutti. 
	  } 
	   } 
	   } 
	 & & \\ 
 

    \theutterance \stepcounter{utterance}  

    & & \multicolumn{4}{p{0.6\linewidth}}{\cellcolor[rgb]{0.9,0.9,0.9}{%
	\makecell[{{p{\linewidth}}}]{% 
	  \tt {\tiny [A$\langle$GM]}  
	 Grazie per la trasparenza. Ecco i miei valori: telefono 6, specchio 0, frutta 2 ciascuno.\\ \tt \\ \tt Dato che entrambi valutiamo molto il telefono, ma io lo valuto 6 contro il tuo 4, l’allocazione efficiente per massimizzare la somma totale è:\\ \tt - Telefono a me (valore totale 6 per me vs 4 per te)\\ \tt - Specchi a te (valgono 0 per me e 2 ciascuno per te)\\ \tt - Frutti possiamo dividerli: io ne prendo almeno 1 (vale 2 per me), tu l’altro (vale 1 per te). Somma: 6 + 4 + 2 + 1 = 13.\\ \tt \\ \tt Se sei d’accordo, posso inviare subito una proposta segreta con: 1 telefono, 1 frutto, e lascio a te 2 specchi e 1 frutto. Fammi sapere se preferisci prendere entrambi i frutti tu (per me va bene rinunciare a 2 punti se questo facilita l’accordo), oppure confermi la divisione 1-1 dei frutti. 
	  } 
	   } 
	   } 
	 & \\ 
 

    \theutterance \stepcounter{utterance}  

    & \multicolumn{4}{p{0.6\linewidth}}{\cellcolor[rgb]{0.8,1,0.9}{%
	\makecell[{{p{\linewidth}}}]{% 
	  \tt {\tiny [A$\rangle$GM]}  
	 La tua proposta massimizza bene la somma. Per me va benissimo: telefono a te, due specchi a me. Sui frutti accetto la divisione 1-1 come suggerito, così otteniamo totale 13.\\ \tt \\ \tt Se sei d’accordo, inviamo entrambi subito le proposte segrete coerenti. Io formulerò la mia per prendere 2 specchi e 1 frutto. Confermi che tu invierai la proposta per 1 telefono e 1 frutto? 
	  } 
	   } 
	   } 
	 & & \\ 
 

    \theutterance \stepcounter{utterance}  

    & & & \multicolumn{2}{p{0.3\linewidth}}{\cellcolor[rgb]{0.95,0.95,0.95}{%
	\makecell[{{p{\linewidth}}}]{% 
	  \tt {\tiny [GM$|$GM]}  
	 continue 
	  } 
	   } 
	   } 
	 & & \\ 
 

    \theutterance \stepcounter{utterance}  

    & & & \multicolumn{2}{p{0.3\linewidth}}{\cellcolor[rgb]{0.95,0.95,0.95}{%
	\makecell[{{p{\linewidth}}}]{% 
	  \tt {\tiny [GM$|$GM]}  
	 La tua proposta massimizza bene la somma. Per me va benissimo: telefono a te, due specchi a me. Sui frutti accetto la divisione 1-1 come suggerito, così otteniamo totale 13.\\ \tt \\ \tt Se sei d’accordo, inviamo entrambi subito le proposte segrete coerenti. Io formulerò la mia per prendere 2 specchi e 1 frutto. Confermi che tu invierai la proposta per 1 telefono e 1 frutto? 
	  } 
	   } 
	   } 
	 & & \\ 
 

    \theutterance \stepcounter{utterance}  

    & & \multicolumn{4}{p{0.6\linewidth}}{\cellcolor[rgb]{0.9,0.9,0.9}{%
	\makecell[{{p{\linewidth}}}]{% 
	  \tt {\tiny [GM$\rangle$B]}  
	 La tua proposta massimizza bene la somma. Per me va benissimo: telefono a te, due specchi a me. Sui frutti accetto la divisione 1-1 come suggerito, così otteniamo totale 13.\\ \tt \\ \tt Se sei d’accordo, inviamo entrambi subito le proposte segrete coerenti. Io formulerò la mia per prendere 2 specchi e 1 frutto. Confermi che tu invierai la proposta per 1 telefono e 1 frutto? 
	  } 
	   } 
	   } 
	 & \\ 
 

    \theutterance \stepcounter{utterance}  

    & & & \multicolumn{4}{p{0.6\linewidth}}{\cellcolor[rgb]{1,0.85,0.72}{%
	\makecell[{{p{\linewidth}}}]{% 
	  \tt {\tiny [GM$\langle$B]}  
	 Confermo. Invio ora la proposta segreta per 1 telefono e 1 frutto.\\ \tt \\ \tt [Proposta: 1 telefono, 1 frutto] 
	  } 
	   } 
	   } 
	  \\ 
 

    \theutterance \stepcounter{utterance}  

    & & & \multicolumn{2}{p{0.3\linewidth}}{\cellcolor[rgb]{0.95,0.95,0.95}{%
	\makecell[{{p{\linewidth}}}]{% 
	  \tt {\tiny [GM$|$GM]}  
	 proposal 
	  } 
	   } 
	   } 
	 & & \\ 
 

    \theutterance \stepcounter{utterance}  

    & & & \multicolumn{2}{p{0.3\linewidth}}{\cellcolor[rgb]{0.95,0.95,0.95}{%
	\makecell[{{p{\linewidth}}}]{% 
	  \tt {\tiny [GM$|$GM]}  
	 [1, 0, 1] 
	  } 
	   } 
	   } 
	 & & \\ 
 

    \theutterance \stepcounter{utterance}  

    & & \multicolumn{4}{p{0.6\linewidth}}{\cellcolor[rgb]{0.9,0.9,0.9}{%
	\makecell[{{p{\linewidth}}}]{% 
	  \tt {\tiny [A$\langle$GM]}  
	 L'altro giocatore ha presentato una proposta segreta. Ora tocca a Lei presentare una sua proposta.\\ \tt  
	  } 
	   } 
	   } 
	 & \\ 
 

    \theutterance \stepcounter{utterance}  

    & \multicolumn{4}{p{0.6\linewidth}}{\cellcolor[rgb]{0.8,1,0.9}{%
	\makecell[{{p{\linewidth}}}]{% 
	  \tt {\tiny [A$\rangle$GM]}  
	 [Proposta: 1 telefono, 1 frutto] 
	  } 
	   } 
	   } 
	 & & \\ 
 

    \theutterance \stepcounter{utterance}  

    & & & \multicolumn{2}{p{0.3\linewidth}}{\cellcolor[rgb]{0.95,0.95,0.95}{%
	\makecell[{{p{\linewidth}}}]{% 
	  \tt {\tiny [GM$|$GM]}  
	 proposal 
	  } 
	   } 
	   } 
	 & & \\ 
 

    \theutterance \stepcounter{utterance}  

    & & & \multicolumn{2}{p{0.3\linewidth}}{\cellcolor[rgb]{0.95,0.95,0.95}{%
	\makecell[{{p{\linewidth}}}]{% 
	  \tt {\tiny [GM$|$GM]}  
	 [1, 0, 1] 
	  } 
	   } 
	   } 
	 & & \\ 
 

    \theutterance \stepcounter{utterance}  

    & & & \multicolumn{2}{p{0.3\linewidth}}{\cellcolor[rgb]{0.95,0.95,0.95}{%
	\makecell[{{p{\linewidth}}}]{% 
	  \tt {\tiny [GM$|$GM]}  
	 [[1, 0, 1], [1, 0, 1]] 
	  } 
	   } 
	   } 
	 & & \\ 
 

\end{supertabular}
}

\end{document}
